\section{Algebraic transfer and an obstructed representative}\label{sec:v7-bypass}

The statements \Fii{} and \Fiii{} impose different requirements on a Hodge
class.  The first asks for a cycle on an abelian variety.  The second asks
for an algebraic correspondence from an abelian variety.  We record the
logical implications precisely, and then give an explicit obstruction to
deducing deformation of a chosen representative from its Chern character.

\subsection{Absolute Hodge classes and algebraic transfer}

Let $Y$ be smooth and projective over $\CC$.  The form of \Fiii{} used here
means that every $\alpha\in\Hdg^p(Y)$ has an expression
\begin{equation}\label{v7:transfer-expression}
 \alpha=\sum_{i=1}^{r}(\Gamma_i)_*\beta_i,
 \qquad
 \beta_i\in\Hdg^{q_i}(A_i),\quad
 \Gamma_i\in\operatorname{CH}^{\dim A_i+p-q_i}(A_i\times Y)_{\QQ},
\end{equation}
where the $A_i$ are complex abelian varieties.  A point is allowed as an
abelian variety of dimension zero.  The action of a correspondence is
\[
 (\Gamma_i)_*\beta_i
   =(\operatorname{pr}_Y)_*
       \bigl(\operatorname{pr}_{A_i}^*\beta_i\cup[\Gamma_i]\bigr).
\]

\begin{proposition}\label{v7:absolute-transfer}
If $Y$ satisfies \Fiii{}, every rational Hodge class on $Y$ is an absolute
Hodge class.  Consequently the universal statement \Fiii{} implies the
absolute Hodge conjecture for all smooth projective complex varieties.
\end{proposition}

\begin{proof}
Use an expression \eqref{v7:transfer-expression}.  Deligne's theorem
\cite[Theorem~2.11]{Del82} makes every $\beta_i$ an absolute Hodge class.
Cycle classes of algebraic correspondences are absolute Hodge, and
pullback, cup product and proper pushforward preserve absolute Hodge
classes.  To see the last assertion directly, spread the finitely many
varieties and correspondences over one subfield of $\CC$ of finite
transcendence degree.  Under every embedding of that field into $\CC$,
comparison and base change identify the displayed correspondence action
with the action of the conjugate algebraic correspondence.  Its source
class remains rational and of Hodge type by Deligne's theorem.  Its image
therefore remains rational and of Hodge type with the indicated Tate
twist.  Finite sums have the same property, so $\alpha$ is absolute Hodge.
\end{proof}

\begin{proposition}\label{v7:equivalence-and-divisors}
The rational Hodge conjecture for all smooth projective complex varieties
is equivalent to the conjunction of the universal statements \Fii{} and
\Fiii{}.  For a fixed $Y$, the following two conditions are equivalent.
\begin{enumerate}
\item Every rational Hodge class on $Y$ is algebraic.
\item Every rational Hodge class on $Y$ admits
  \eqref{v7:transfer-expression} with each $\beta_i$ in the algebra
  generated by divisor classes on $A_i$.
\end{enumerate}
The second condition has the same force if ``in the algebra generated by
divisor classes'' is replaced by ``algebraic.''
\end{proposition}

\begin{proof}
Assume \Fii{} and \Fiii{}, and write a Hodge class as in
\eqref{v7:transfer-expression}.  By \Fii{} choose cycles $z_i$ with
$[z_i]=\beta_i$.  Compatibility of the cycle class map with intersection
and proper pushforward gives
\[
 \alpha=
 \left[\sum_i(\operatorname{pr}_Y)_*
    \bigl(\operatorname{pr}_{A_i}^*z_i\cdot\Gamma_i\bigr)\right].
\]
This proves the rational Hodge conjecture.  Conversely that conjecture
includes \Fii{}.  If $\alpha=[z]$ on $Y$, take the source to be a point,
$\beta=1$ and the correspondence to be $z$.  This proves \Fiii{}.

For the fixed-variety assertion, a polynomial in divisor classes is the
class of the same polynomial in divisor cycles.  The preceding
pushforward formula proves that condition~2 implies condition~1, and
also proves the variant with algebraic source classes.  For the converse
one can again use a point.  If a positive-dimensional source is required,
fix an elliptic curve $A$, a point $a\in A$ and a cycle $z$ representing
$\alpha$.  Then $\beta=[a]\in\Hdg^1(A)$ is a divisor class and the
correspondence $A\times z$ satisfies
$(A\times z)_*[a]=[z]=\alpha$.  Thus the divisor-source formulation
already has exactly the force of the rational Hodge conjecture on $Y$.
\end{proof}

For each abelian variety $A$, statement \Fiii{} itself holds by the
diagonal correspondence, taking the source class to be the given Hodge
class on $A$.  Thus this special case supplies no algebraic
representative of that source class.  Proposition~\ref{v7:absolute-transfer}
also supplies absoluteness without an algebraic representative.  These
are the precise conclusions available from Deligne's theorem and
algebraic transfer.

\subsection{A polystable bundle whose full character stays Hodge}

For a perfect complex $E$ on a smooth projective variety $S$, its
first-order obstruction along a deformation with Kodaira--Spencer class
$\kappa\in H^1(S,T_S)$ is
\begin{equation}\label{v7:atiyah-obstruction}
 o_{\kappa}(E)=\kappa\lrcorner\operatorname{At}(E)
       \in\Ext^2(E,E).
\end{equation}
We use the common normalization in which the Atiyah class of a line
bundle is its first Chern class.  Changing the convention changes the
formula by a common nonzero scalar or sign and preserves its vanishing.
For a line bundle $M$, the obstruction is therefore
$\kappa\lrcorner c_1(M)\in H^2(S,\mathcal O_S)$.
This formula can also be obtained from its Atiyah extension
\[
 0\longrightarrow\mathcal O_S
  \longrightarrow\mathcal D^1(M)
  \longrightarrow T_S\longrightarrow0.
\]
The connecting map on $H^1(T_S)$ is contraction with the extension
class $c_1(M)$ and measures the failure of the transition functions of
$M$ to lift over the prescribed deformation of $S$.

\begin{theorem}\label{v7:k3-obstructed-bundle}
There exist a smooth quartic K3 surface $S$, its hyperplane polarization
$H$, an $H$-polystable rank-two vector bundle $E$, and an $H$-polarized
first-order deformation $S_{\varepsilon}$ over
$\CC[\varepsilon]/(\varepsilon^2)$ with the following properties.
\begin{enumerate}
\item If $\eta\in H^4(S,\QQ)$ is the class of a point, then
  $\ch(E)=2-36\eta$.  Every component of this character extends as a
  flat Hodge class over the whole local polarized K3 deformation space.
\item The bundle $E$, viewed as a perfect complex, has no deformation
  to $S_{\varepsilon}$.
\item There are identifications
  \[
   \Ext^2(E,E)\cong\CC\oplus\CC,
   \qquad
   \sigma_E(a,b)=a+b,
  \]
  where $\sigma_E$ is the semiregularity map.  Its one-dimensional
  kernel contains the nonzero obstruction to this first-order lift.
\end{enumerate}
\end{theorem}

\begin{proof}
One explicit surface is the Fermat quartic
\[
 S=\{x_0^4+x_1^4+x_2^4+x_3^4=0\}\subset\mathbb P^3_{\CC}.
\]
Its partial derivatives have no common projective zero, so it is smooth.
Choose $a,b\in\CC$ with $a^4=b^4=-1$.  The equations
$x_0=ax_1$ and $x_2=bx_3$ define a line $\ell\subset S$.
Adjunction gives $K_S\cong\mathcal O_S$, and the hypersurface sequence
gives $H^1(S,\mathcal O_S)=0$, so $S$ is a K3 surface.
Write $H=c_1(\mathcal O_S(1))$.  Its intersection numbers are
\[
 H^2=4,\qquad H\cdot\ell=1,\qquad\ell^2=-2.
\]
The last equality follows by adjunction on the rational curve $\ell$.
Put
\begin{equation}\label{v7:k3-line-class}
 l=4[\ell]-H,\qquad L=\mathcal O_S(4\ell-H),\qquad
 E=L\oplus L^{-1}.
\end{equation}
Then
\begin{equation}\label{v7:k3-intersections}
 H\cdot l=0,\qquad l^2=-36\eta.
\end{equation}
In particular $l\ne0$, and it is independent of $H$.
Every line bundle is slope-stable.  The two summands of $E$ both have
$H$-slope zero, so $E$ is $H$-polystable.  No assertion about the Picard
number of $S$ is needed.

On a surface the Chern character has only three components, and
\[
 \ch(E)=\exp(l)+\exp(-l)=2+l^2=2-36\eta.
\]
The unit in $H^0$ and the normalized point class in $H^4$ are flat
Hodge classes in every local smooth deformation.  This proves the
first assertion.

We choose a polarized tangent direction that does not preserve $l$.
Contraction followed by cup product is the perfect Serre-duality pairing
\[
 H^1(S,T_S)\times H^1(S,\Omega_S^1)
       \longrightarrow H^2(S,\mathcal O_S)\cong\CC,
\]
using $K_S\cong\mathcal O_S$.  Since $H$ and $l$ are independent, there
is a $\kappa\in H^1(S,T_S)$ such that
\begin{equation}\label{v7:k3-tangent}
 \kappa\lrcorner H=0,\qquad
 u:=\kappa\lrcorner l\ne0.
\end{equation}
The first equation says that this is a tangent direction to polarized
deformations.  Local Torelli for K3 surfaces \cite[Theorem~5.7]{LP80} identifies their smooth
local polarized deformation space with an open subset of the period
quadric in $H^{\perp}$.  The extra condition that $l$ stay of type
$(1,1)$ is the additional period hyperplane $l^{\perp}$.
It is a proper hyperplane here.  Indeed $l$ itself represents a tangent
vector to the polarized period quadric at $H^{2,0}(S)$, and the derivative
of the defining equation of $l^{\perp}$ on that vector is $l^2=-36$.
Thus directions satisfying \eqref{v7:k3-tangent} occur in actual local
polarized deformations.  Restrict one to the first infinitesimal
neighbourhood of its central point to obtain $S_{\varepsilon}$.

The Atiyah class of a direct sum is block diagonal, and
$\operatorname{At}(L^{-1})=-\operatorname{At}(L)$.  Consequently
\eqref{v7:atiyah-obstruction} gives
\begin{equation}\label{v7:k3-obstruction-matrix}
 o_{\kappa}(E)=
 \begin{pmatrix}u&0\\0&-u\end{pmatrix}\ne0.
\end{equation}
Nonvanishing is meant in the direct-sum decomposition of
$\Ext^2(E,E)$, so deformation of the splitting is not assumed.
The obstruction class rules out a lift of the object $E$ itself,
including a lift as a perfect complex.

For completeness we compute the whole obstruction space.  Each diagonal
summand is
$\Ext^2(L,L)=\Ext^2(L^{-1},L^{-1})=H^2(S,\mathcal O_S)\cong\CC$.
Moreover $H^0(S,L^2)=H^0(S,L^{-2})=0$.  A nonzero section of either
bundle would define a nonzero effective divisor of $H$-degree zero,
contradicting ampleness.  A section without zeros would trivialize that
bundle, contradicting $l^2\ne0$.  Serre duality therefore gives
\[
 \Ext^2(L,L^{-1})=H^2(S,L^{-2})=0,
 \qquad
 \Ext^2(L^{-1},L)=H^2(S,L^2)=0.
\]
On a surface the semiregularity target has only its degree-zero Atiyah
component, because $H^{q+2}(S,\Omega_S^q)=0$ for $q>0$.
Thus $\sigma_E$ is the ordinary trace to $H^2(S,\mathcal O_S)$,
which is $(a,b)\mapsto a+b$ in the displayed decomposition.
Its kernel is the anti-diagonal, containing the nonzero class $(u,-u)$
of \eqref{v7:k3-obstruction-matrix}.  Each line-bundle summand separately
has an isomorphism as its semiregularity map.
\end{proof}

All components of the character in Theorem~\ref{v7:k3-obstructed-bundle}
are already algebraic.  The failure concerns the lift of this particular
polystable representative.  The example proves that preserving an entire
Chern character, including the determinant, does not force a chosen
perfect complex to deform.  It also proves that injective semiregularity
of the summands does not imply injective semiregularity of their sum.

\subsection{The remaining requirement for a pure character}

For a general perfect complex, the Atiyah--Kodaira--Spencer obstruction is described in \cite{HT10}; the components of semiregularity have
the form
\[
 \sigma_{E,q}(a)=\frac{1}{q!}
    \operatorname{tr}\bigl(a\operatorname{At}(E)^q\bigr)
       \in H^{q+2}(Y,\Omega_Y^q).
\]
With the signs and normalization of \eqref{v7:atiyah-obstruction},
compatibility with first-order variation gives
\[
 \sigma_{E,q}(o_{\kappa}(E))
       =\kappa\lrcorner\ch_{q+1}(E).
\]
If all components of $\ch(E)$ stay Hodge, these images vanish.
Theorem~\ref{v7:k3-obstructed-bundle} exhibits a nonzero obstruction in
their common kernel.

An Adams-operation construction of a pure cohomological character
$\ch(E)=N\xi_c$ establishes exactly the character condition in this
calculation.  To propagate that representative one still needs an
argument that its obstruction vanishes, and corresponding control at
higher orders.  Injective semiregularity is one sufficient additional
condition.  Equality in $K_0$ leaves the self-extension groups and the
Atiyah class of a representative undetermined.  Likewise, additivity
of the Chern character in a distinguished triangle gives no vanishing
of the obstruction of its cone.  These requirements persist after
the character has been projected to a single degree.

\subsection{Rational domination and total-space hypotheses}

\begin{proposition}\label{v7:rational-varieties}
The universal statement \Fiii{} for smooth projective rational varieties
is equivalent to the universal statement \Fiii{} for all smooth
projective complex varieties.  The analogous equivalence holds for the
rational Hodge conjecture.  Consequently an extension of descent of
\Fiii{} from surjective morphisms to dominant rational maps, even just
generically finite ones, would imply universal \Fiii{}.
\end{proposition}

\begin{proof}
Embed a given smooth projective variety $Y$ of dimension $m$ in
$\mathbb P^N$ with codimension $c=N-m\ge2$, and put
$W=\operatorname{Bl}_Y\mathbb P^N$.  This is smooth, projective and
rational.  Write $j:D\hookrightarrow W$ for its exceptional divisor,
$p:D\to Y$ for the projective bundle of relative dimension $c-1$, and
$h=c_1(\mathcal O_D(1))$.  The algebraic blow-up correspondences
\cite[\S9]{Man68} include
\[
 a(\alpha)=j_*p^*\alpha,\qquad
 b(x)=-p_*\bigl(h^{c-2}j^*x\bigr).
\]
Here $a:H^k(Y)\to H^{k+2}(W)$ is induced by a cycle in
$\operatorname{CH}^{m+1}(Y\times W)$, and
$b:H^{k+2}(W)\to H^k(Y)$ by a cycle in
$\operatorname{CH}^{N-1}(W\times Y)$.
The normal bundle of $D$ in $W$ is $\mathcal O_D(-1)$, so the
self-intersection formula and $p_*h^{c-1}=1$ give
\[
 b(a(\alpha))
   =-p_*\bigl(h^{c-2}(-h)p^*\alpha\bigr)=\alpha.
\]
Thus $a$ embeds $H^{2q}(Y)(q)$ in $H^{2q+2}(W)(q+1)$ with an
algebraic retraction.  If $W$ satisfies \Fiii{}, apply it to $a(\alpha)$
and compose each resulting correspondence with $b$.  This gives
\Fiii{} for $Y$.  If the Hodge conjecture holds for $W$, represent
$a(\alpha)$ by an algebraic cycle and apply $b$ to obtain an algebraic
representative of $\alpha$.  The reverse universal implications follow
by restricting the class of varieties.

Finally $(\mathbb P^1)^N$ is birational to $\mathbb P^N$, and hence has a
dominant generically finite rational map to $W$.  Its cohomology is
generated by divisors.  A descent principle for \Fiii{} under these
rational maps would give \Fiii{} on every such $W$ and therefore on
every $Y$ by the preceding retraction.
\end{proof}

\begin{proposition}\label{v7:total-space-hypothesis}
Let $C$ be a smooth connected projective curve, let $\mathcal X\to C$
and $\mathcal Z\to C$ be projective abelian schemes of relative
dimensions four and six, and put
\[
 f:\mathcal Y=\mathcal X\times_C\mathcal X\times_C\mathcal Z\to C.
\]
Let $\xi\in H^0(C,R^4f_*\QQ)$ be a flat Hodge class, algebraic on one
fibre.  Either of the following additional hypotheses suffices to make
$\xi_s$ algebraic at every $s\in C$.
\begin{enumerate}
\item At one point $c$, every monodromy-invariant rational Hodge class
  in $H^4(\mathcal Y_c,\QQ)$ is the restriction of a rational
  algebraic class of codimension two on $\mathcal Y$.
\item The Lefschetz standard conjecture $B(\mathcal Y)$ holds for the
  total space of this abelian scheme.
\end{enumerate}
The first hypothesis includes an algebraic lift of the missing class
$\xi_c$.  The second is a hypothesis about the mixed total space, of
dimension fifteen.  Statements $B(\mathcal X)$ and $B(\mathcal Z)$
concerning the separate total spaces do not supply either hypothesis
in this criterion.
\end{proposition}

\begin{proof}
For the first assertion, $\xi_c$ is monodromy invariant because $\xi$ is
a global flat section.  The hypothesis gives
$z\in\operatorname{CH}^2(\mathcal Y)_{\QQ}$ with
$i_c^*[z]=\xi_c$, where $i_s:\mathcal Y_s\hookrightarrow\mathcal Y$
denotes the fibre inclusion.  The restrictions $i_s^*[z]$ form a flat
section of $R^4f_*\QQ$.  Since $C$ is connected, equality at $c$
implies equality with $\xi$ everywhere.  Refined pullback by the
codimension-one regular embedding $i_s$ preserves the codimension of
cycles, so $i_s^!z\in\operatorname{CH}^2(\mathcal Y_s)_{\QQ}$
represents $\xi_s$.

For the second assertion, apply the algebraic propagation theorem for
an abelian scheme over a smooth projective base, in the formulation
of \cite[\S5, Proposition~1]{Milne20}; see also the published propagation theorem \cite{Abd94}.  Its total-space hypothesis is
$B(\mathcal Y)$ and its conclusion applies to the given flat section,
which is algebraic on one fibre.

The first assertion shows directly what a proposed total-space cycle
restriction hypothesis must contain.  In the second assertion the
scheme to which the propagation theorem is applied is $f$.  Its total
space is $\mathcal Y$, rather than either separate family.
Algebraic Lefschetz inverses on $\mathcal X$, on $\mathcal Z$ and on
their ordinary products are operators on those spaces.  Passing to
the fibre product requires additional correspondences controlling
its cohomology, including the mixed summand containing $\xi$.
Neither of the stated separate standard conjectures asserts such a
correspondence or an algebraic restriction lift of $\xi_c$.
Thus the criterion still requires the indicated hypothesis on the
mixed family.
\end{proof}
